\section{Related work}\label{sec:related}

\textbf{Evaluating prognostics through decisions.} Remaining useful life (RUL) estimates are usually scored at times chosen by the evaluator, with accuracy metrics or asymmetric scores that penalise late predictions~\cite{saxena2008,saxena2010,ramasso2014}. A line of work in this journal evaluates prognostics through the maintenance decisions they trigger. de~Pater et al.~\cite{depater2022} schedule aircraft maintenance from alarms raised on RUL prognostics and protect against prognostic error with a safety factor; Kamariotis et al.~\cite{kamariotis2024} propose a metric $M$, the relative excess long-run cost of a policy over perfect prognostics, and tune probability-threshold policies with it; dynamic predictive-maintenance policies act on predicted failure probabilities~\cite{nguyen2019,mitici2023}. Age replacement is the benchmark that any condition-based policy must beat~\cite{barlow1960,dejonge2020}. None of these gives a finite-sample guarantee on the notice an alarm leaves.

\textbf{Conformal prediction for RUL.} Conformal prediction turns any point predictor into prediction sets with finite-sample, distribution-free coverage under exchangeability~\cite{vovk2005,lei2018,angelopoulos2023}. Javanmardi and H\"ullermeier~\cite{javanmardi2023} conformalised gradient boosting and convolutional networks to obtain RUL intervals on C-MAPSS; Robinson~\cite{robinson2026} used asymmetric conformalised quantile regression~\cite{romano2019} to make the lower RUL bound conservative; Wang et al.~\cite{wang2025robustuq} combined conformal prediction with imputation of missing sensor data and noted that exchangeability is a concern for degradation data. In all of these the calibrated object is the interval at a given cycle. Cand\`es et al.~\cite{candes2023} give conformal lower bounds on survival times under censoring, at a fixed covariate value and under assumptions on the censoring mechanism. An alarm, however, is a stopping rule, evaluated at the cycle at which the bound first crosses a level; per-cycle coverage does not imply coverage at that time, and Section~\ref{sec:res-valid} shows that the difference is large.

\textbf{Risk control, cross-fitting and shift.} Conformal risk control~\cite{angelopoulos2024} calibrates a scalar parameter so that the expected value of a monotone loss is bounded, and nested conformal prediction~\cite{gupta2022} expresses conformal procedures through nested families of sets. Our construction is an instance of the nested view, with the warning level as the nesting parameter and failure before the planned replacement as the loss. What is specific is the closed-form critical level of a debounced trigger (Lemma~\ref{lem:crit}), the pathwise censoring bound that makes calibration possible with replaced units (Theorem~\ref{thm:cens}), the decision-aware choice that keeps the guarantee (Corollary~\ref{cor:da}), and the certification floor (Proposition~\ref{prop:cost}). Cross-fitting follows the jackknife+/CV+ construction of Barber et al.~\cite{barber2021}. When calibration and new units are not exchangeable, weighted and non-exchangeable conformal methods bound the coverage gap~\cite{tibshirani2019,barber2023}, but they need a model of the shift; online methods instead track the target error rate under arbitrary shift by updating a level from observed errors~\cite{gibbs2021,angelopoulos2023pid}. We use the latter idea on the critical level (Theorem~\ref{thm:online}), where it has a natural feedback signal: whether each unit failed before its replacement. We also measure how badly exchangeability fails across battery batches and N-CMAPSS subsets, and what pilot units cost.

\textbf{Relation to our earlier evaluation.} This article supersedes an evaluation study of the same alarm recipe on the same data~\cite{prior}, which separated prediction accuracy, alarm timing, maintenance cost and transfer and proposed no method. Its numbers are cited as prior results; every number presented as new was produced by the code released with this article.

\section{Problem formulation}\label{sec:problem}

\textbf{Units, alarms and notice.} A unit observed at rows $j=1,\dots,n_u$ (cycles $t_1<\dots<t_{n_u}$) has life $L$ and remaining life $r_j=L-t_j\ge1$, decreasing in $j$. A causal signal $s_j$, here a smoothed RUL estimate, is computed from rows $1,\dots,j$. A \emph{threshold trigger} with debounce $k$ and warning level $H$ raises an alarm at the first row at which the signal has been at or below $H$ for $k$ consecutive rows. Writing $m_j=\max(s_{j-k+1},\dots,s_j)$ for $j\ge k$ ($m_j=+\infty$ otherwise) and $\rho_j=\min_{i\le j}m_i$ for the \emph{debounced running minimum}, the alarm row is
\[
\tau(H)=\min\{j:\rho_j\le H\}\qquad(\tau=\infty\text{ if there is no such row}),
\]
and the \emph{notice} is $N(H)=r_{\tau(H)}$ ($N=0$ without an alarm).

\textbf{Maintenance consequence.} An alarm schedules a replacement $\ell$ cycles later (the \emph{lead time}). If $N(H)>\ell$ the planned replacement happens at $t_{\tau(H)}+\ell$ at cost~1; otherwise, or without an alarm, the unit fails at $L$ at cost $\rho>1$. Units are renewed, so by the renewal--reward theorem the long-run cost per cycle is $E[\text{cost}]/E[\text{use}]$, estimated by total cost over total cycles of use and reported relative to perfect foresight (replacement at $L-1$, cost~1)~\cite{barlow1960}. The reliability requirement we target is
\[
P\bigl(N(H)\le\ell\bigr)\le\alpha ,
\]
the probability that the unit fails before the replacement its alarm scheduled.

\textbf{Window score.} The earlier study scored the first alarm by an asymmetric window: with reported estimate $\hat Y$ and error $e=\hat Y-R$, $R=N(H)$, the alarm is correct if $a\le e\le b$, with $[a,b]=[-13,+2]$ cycles for the engines and $[-39,+6]$ cycles for the battery cells.
